\clearpage
\fancyhead[L]{\small PALLAS. JUNO.}
\hypertarget{native-274-lower}{}
\begin{center}
PALLAS. JUNO.

\rule{0.6\linewidth}{0.4pt}

Monatliche Correspondenz. Vol. XIII. Pp. 313–315. March 1806.

\rule{0.6\linewidth}{0.4pt}
\end{center}

\par\indent So far the only observations of \emph{Pallas} that have come to our attention are the three made at the Seeberg Observatory and reported in the previous issue, \href{https://www.digitale-sammlungen.de/en/view/bsb10538605?page=197}{p. 189}. We have yet to hear of any observations of \emph{Juno}. We therefore hasten to share the following news with our readers.

\par\indent Dr.\ \textsc{Gauss} first saw \emph{Pallas} again on 13 February. He made only this one observation, solely to pick out the planet by its movement on the next clear evening from among the other nearby stars, all of which he had to note as well because of the uncertainty. He therefore did not send it to us at all, considering it very imprecise. On the following evening, 14 February, it was now easy for him to recognize Pallas, which had moved and looked like a star of the 7\textsuperscript{th} to 8\textsuperscript{th} magnitude. That evening he compared it with 54 Eridani, and on the following evenings with stars from the Hist.\ cél. Here are the observations [the same as on \hyperlink{native-272}{p. 267} of this volume].

\par\indent According to these observations, Dr.\ \textsc{Gauss}’s ephemeris in the M.\,C., vol.\ XI, \href{https://www.digitale-sammlungen.de/en/view/bsb10538603?page=388}{p. 376} [\hyperlink{pallas180504.1}{p. 263} of this volume] gives this planet’s right ascension 27 minutes too high and its declination 5 minutes too high. “It is strange,” Dr.\ \textsc{Gauss} remarks, “that the first discrepancy is exactly in the opposite direction to that indicated by your observations last November (see the previous issue, \href{https://www.digitale-sammlungen.de/en/view/bsb10538605?page=197}{p. 189}). The size of the error now shows, I think, the effect of the perturbations.”

\par\indent He first saw \emph{Juno} again on 16 February and on 17 February observed it again \emph{as such}. Here are his observations [the same as on \hyperlink{native-272}{p. 267} of this volume].

\par\indent On this evidence, Dr.\ \textsc{Gauss}’s ephemeris (M.\,C., vol.\ XI, \href{https://www.digitale-sammlungen.de/en/view/bsb10538603?page=491}{p. 477} [\hyperlink{native-270}{p. 265} of this volume]) would give a right ascension one or one and a half minutes too high and a declination one minute too low. After an interval of exactly one year, this is a quite insignificant discrepancy. We can also conclude from it that Dr.\ \textsc{Gauss}’s last observation, on 20 February 1805, on which his fifth set of ele%
\SourceLeaf{275}%
ments was based, must have been exceptionally good. Dr.~\textsc{Gauss} writes on this occasion, “Juno is very faint. I estimate it to be about as bright as a star of the eleventh magnitude. In clear air it can still be observed well with the ring micrometer, but it will scarcely be possible with a transit instrument and a circle or quadrant if illumination is needed. I usually find these observations very troublesome, since I have only a reflecting telescope made by \textsc{Short} for this work. Its field of view is very small, so most of the time there are no suitable stars of accurately known position close enough to pass through the narrow field along with the planet. There is also the often substantial difficulty of finding such small planets in the first place when there are no known stars nearby to serve as guide stars.”

\par\indent We add one more piece of news that is sure to be most welcome to our readers interested in astronomy. Dr.~\textsc{Gauss} tells us that he is now working chiefly on several particular topics in his method of determining planetary orbits, and that he hopes soon to be able to devote himself in earnest to working out, completing and publishing a separate work on the subject.

\begin{flushright}
[\textsc{von Zach}.]
\end{flushright}
